% ============================================================
% Corpo principal do guia. Incluído a partir de main.tex.
% Headings em nível de \chapter (classe book).
% ============================================================

\chapter{O ponto de partida: a regressão linear}
% ============================================================

A regressão linear é, de longe, o método mais usado em econometria aplicada. E há boas
razões para isso: ela é simples de estimar, fácil de interpretar e, em muitos contextos,
fornece uma aproximação razoável da relação entre variáveis. Quando queremos resumir como
uma variável $X$ se relaciona com um resultado $Y$, o modelo linear diz:

\begin{equation}
Y = \alpha + \beta X + \varepsilon,
\label{eq:linear}
\end{equation}

\noindent
onde $\alpha$ é o intercepto, $\beta$ é o coeficiente de inclinação e $\varepsilon$
captura tudo o mais que afeta $Y$ mas não está incluído no modelo.

O que $\beta$ significa? Ele mede o efeito médio de uma unidade adicional de $X$ sobre
$Y$ — e, criticamente, esse efeito é \emph{o mesmo} para qualquer nível de $X$. Se
$\beta = 1{,}5$, cada unidade extra de $X$ está associada, em média, a um aumento de
1,5 unidades em $Y$, independentemente de $X$ valer 2 ou 200.

Essa é a essência do modelo linear: \textbf{um único coeficiente resume a relação
entre $X$ e $Y$ em todos os pontos da distribuição}.

\begin{exemplobox}[Exemplo: idade e salário]
Suponha que $Y$ é o salário e $X$ é a idade. Uma regressão linear estimaria que cada
ano adicional de idade aumenta o salário por uma quantidade fixa. Mas pense no que
acontece na vida real: no início da carreira, ganhos salariais costumam ser rápidos;
na meia-idade, os salários se estabilizam; e, perto da aposentadoria, podem até cair.
A relação verdadeira tem a forma de uma curva de sino, não de uma reta. Uma regressão
linear aplicada a esses dados produziria uma reta que nenhum ponto da curva real
representa bem.
\end{exemplobox}

Isso não significa que a regressão linear seja inútil. Em muitos casos, ela captura
a maior parte da variação relevante dos dados, e sua simplicidade facilita enormemente
a comunicação dos resultados. O ponto é que a linearidade é uma \emph{hipótese}, não
um fato. Quando essa hipótese é razoável, o modelo linear é uma ferramenta poderosa.
Quando não é, os resultados podem enganar.

% ============================================================
\chapter{Quando a linearidade vira uma restrição forte demais}
% ============================================================

O problema com a linearidade não é apenas estético. Impor uma forma funcional errada
pode distorcer as conclusões de maneira importante. Vejamos alguns padrões que um
modelo linear captura mal.

\textbf{Retornos decrescentes.} Em muitos contextos, os primeiros anos de educação
têm um impacto enorme sobre o salário; anos adicionais de escolaridade já têm um
retorno menor. Isso é uma relação côncava: o efeito marginal de educação sobre salário
diminui à medida que o nível de educação aumenta. Um modelo linear forçaria um efeito
constante, provavelmente superestimando o retorno da educação em altos níveis e
subestimando-o em baixos.

\textbf{Limiares (\emph{thresholds}).} Em estudos de saúde ambiental, é comum que
a poluição cause pouco dano até um determinado nível crítico, a partir do qual os
efeitos sobre a saúde disparam. Um modelo linear forçaria uma relação proporcional em
toda a faixa de poluição, ignorando essa descontinuidade.

\textbf{Platôs.} Gastos públicos em educação podem melhorar indicadores de
aprendizagem até um certo ponto, mas a partir daí o efeito se esgota. A relação tem
um platô. Uma regressão linear subestimaria o impacto dos primeiros reais gastos e
superestimaria o efeito de gastos já elevados.

\textbf{Formato de U ou U invertido.} A relação entre experiência no trabalho e
produtividade pode crescer no início, atingir um máximo e depois declinar. Ou a
relação entre desigualdade e crescimento econômico pode ter essa forma. Em ambos os
casos, um único coeficiente linear é uma caricatura.

\textbf{Efeitos que variam ao longo da distribuição.} Mesmo quando a relação é
aproximadamente monótona, o tamanho do efeito pode ser diferente para valores
baixos, médios e altos de $X$. Uma regressão linear calcula uma média ponderada desses
efeitos, mas pode mascarar heterogeneidade relevante.

\begin{destacebox}[A questão prática]
Em todos esses casos, o pesquisador que impõe linearidade sem questionar pode chegar a
conclusões quantitativamente erradas e, às vezes, qualitativamente enganosas. O antídoto
é: antes de impor uma forma funcional, olhe os dados. Faça gráficos de dispersão, explore
relações não paramétricas e só imponha estrutura quando houver razão teórica ou empírica
para isso.
\end{destacebox}

Essa reflexão prepara o terreno para os métodos não paramétricos. Se não queremos
impor uma forma funcional, precisamos de ferramentas que deixem os dados revelarem
a relação entre $X$ e $Y$ por conta própria.

% ============================================================
\chapter{A ideia dos métodos não paramétricos}
% ============================================================

A ideia central dos métodos não paramétricos é simples: em vez de impor que a
relação entre $X$ e $Y$ tem uma forma específica, deixamos que essa forma emerja
dos dados.

Em termos formais, em vez de estimar o modelo:

\begin{equation}
Y = \alpha + \beta X + \varepsilon,
\end{equation}

\noindent
queremos estimar:

\begin{equation}
Y = g(X) + \varepsilon,
\label{eq:nonparam}
\end{equation}

\noindent
onde $g(\cdot)$ é uma função \emph{desconhecida}. Não impomos que $g$ seja uma reta,
uma parábola ou qualquer outra forma pré-definida. Queremos estimar $g$ diretamente
dos dados.

\textbf{O que ganhos com isso?} A flexibilidade para descobrir se a relação é
crescente, decrescente, côncava, convexa, com platô, com salto, ou qualquer outra
forma que os dados sugiram. Métodos não paramétricos são excelentes ferramentas
exploratórias: permitem enxergar padrões que uma especificação linear esconderia.

\textbf{O que perdemos?} Liberdade tem um custo.

\begin{itemize}
  \item \textbf{Precisão.} Quanto mais flexível o modelo, mais dados são necessários
  para estimar cada parte da função com precisão. Com poucos dados, estimativas muito
  flexíveis são ruidosas.

  \item \textbf{Viés e variância.} Em estatística, há um dilema clássico: modelos
  mais simples (como o linear) têm viés potencialmente alto se a relação verdadeira
  não for linear, mas têm variância baixa. Modelos muito flexíveis têm viés baixo,
  mas variância alta. Métodos não paramétricos precisam navegar esse dilema.

  \item \textbf{\emph{Overfitting}.} Um modelo excessivamente flexível pode se ajustar
  ao ruído presente na amostra, em vez de capturar o padrão real. O resultado parece
  perfeito nos dados de treinamento, mas generaliza mal.

  \item \textbf{Maldição da dimensionalidade.} Quando há muitas covariáveis $X_1, X_2,
  \ldots, X_k$, estimar $g(X_1, X_2, \ldots, X_k)$ de forma não paramétrica exige
  quantidades enormes de dados. O espaço de covariáveis cresce exponencialmente com
  o número de variáveis, e os dados ficam cada vez mais ``espalhados'' nesse espaço.
\end{itemize}

\begin{destacebox}[Resumo intuitivo do dilema]
Pense em um artesão que pode usar régua ou mão livre para desenhar uma curva. Com a
régua, o traço é preciso, mas só serve para retas. Sem a régua, pode desenhar qualquer
forma — mas pequenos tremores da mão introduzem imperfeições. Métodos não paramétricos
são como desenhar à mão livre: mais expressivos, mas mais sensíveis ao ruído. O
desafio é controlar o quanto a mão ``treme'' sem perder a capacidade de desenhar
curvas.
\end{destacebox}

Na prática, métodos não paramétricos são mais úteis quando: (i) há quantidade
razoável de dados, (ii) o número de covariáveis é pequeno (idealmente uma ou duas),
e (iii) o pesquisador quer explorar a forma da relação antes de impor estrutura.

% ============================================================
\chapter{Estimação local: a intuição por trás da regressão linear local}
% ============================================================

Como, na prática, estimamos uma função $g(X)$ sem impor sua forma? A ideia mais
intuitiva é a \textbf{estimação local}: se queremos saber como $Y$ se comporta perto
de um ponto específico $x_0$, faz sentido olhar principalmente para as observações
que estão \emph{próximas} de $x_0$.

Observações muito distantes de $x_0$ podem ser pouco informativas sobre o
comportamento local da função. Se a relação é não linear, o que acontece em $X = 2$
pode dizer pouco sobre o que acontece em $X = 10$. Portanto, ao estimar o valor de
$g$ próximo de $x_0$, podemos dar mais peso às observações vizinhas e menos às mais
distantes.

\section*{Da média local à regressão linear local}

A forma mais simples de estimação local seria calcular a média de $Y$ para as
observações com $X$ próximo de $x_0$. Isso é chamado de \emph{Nadaraya-Watson
estimator} e produz, para cada ponto $x_0$, uma média ponderada de $Y$.

A \textbf{regressão linear local} (\emph{local linear regression}) vai um passo além:
em vez de apenas calcular uma média ponderada, estima uma \emph{reta} dentro de cada
vizinhança local. Formalmente, para cada ponto $x_0$, minimizamos:

\begin{equation}
\sum_{i=1}^{n} K\!\left(\frac{X_i - x_0}{h}\right) \bigl(Y_i - a - b(X_i - x_0)\bigr)^2,
\label{eq:llr}
\end{equation}

\noindent
onde $K(\cdot)$ é o \emph{kernel} (que define os pesos), $h$ é o \emph{bandwidth}
(que define o tamanho da vizinhança) e $(a, b)$ são os parâmetros locais. O
estimador de $g(x_0)$ é o valor ajustado $\hat{a}$, que corresponde à reta local
avaliada no ponto $x_0$.

\section*{O papel do kernel e do bandwidth}

\textbf{Kernel} ($K$): define como os pesos diminuem com a distância até $x_0$.
Kernels populares atribuem peso máximo para observações em $x_0$ e decrescem suavemente
até zero nas bordas da vizinhança. O kernel mais comum é o gaussiano, mas outros
(como Epanechnikov ou triangular) também são usados.

\textbf{Bandwidth} ($h$): controla o tamanho da vizinhança. Este é o parâmetro
mais importante da estimação:

\begin{itemize}
  \item \textbf{Bandwidth pequeno}: a vizinhança é estreita, usamos poucas observações,
  a estimativa é muito flexível (segue os dados de perto), mas pode ser bastante
  ruidosa.
  \item \textbf{Bandwidth grande}: a vizinhança é larga, usamos muitas observações,
  a estimativa é mais suave e estável, mas pode não capturar variações locais
  importantes.
\end{itemize}

\begin{exemplobox}[Voltando ao exemplo de idade e salário]
Imagine uma amostra de trabalhadores com suas idades e salários. Uma regressão linear
produziria uma reta: talvez um efeito positivo de \$500 por ano adicional de idade.
Uma regressão linear local revelaria algo diferente: crescimento salarial rápido entre
20 e 35 anos (trabalhadores iniciantes ganhando experiência), estabilização entre 35
e 50 anos, e eventual declínio depois dos 55. O resultado é uma curva suave que
conta uma história muito mais rica do que a reta linear.
\end{exemplobox}

\section*{Por que regressão linear local e não apenas médias locais?}

A regressão linear local tem uma vantagem técnica importante sobre estimadores de
médias ponderadas: ela funciona melhor nas \textbf{bordas da distribuição de $X$}.

Quando $x_0$ está perto do extremo da distribuição de $X$ (por exemplo, idades muito
jovens ou muito velhas na amostra), o estimador de média local sofre de um viés de
fronteira, pois as observações disponíveis ficam todas de um lado de $x_0$. A
regressão linear local incorpora a inclinação local, o que reduz substancialmente
esse problema.

\begin{notatecnica}[Nota técnica: escolha do bandwidth]
Na prática, $h$ é escolhido por \emph{cross-validation}: o pesquisador testa vários
valores de $h$ e escolhe aquele que minimiza o erro de previsão fora da amostra.
Regras de bolso (\emph{rules of thumb}) baseadas na distribuição normal também são
comuns. O pacote \texttt{nprobust} em R e Stata oferece métodos automáticos de
seleção de bandwidth com propriedades ótimas.
\end{notatecnica}

% ============================================================
\chapter{Até aqui, o problema era forma funcional --- agora entra a endogeneidade}
% ============================================================

Até agora, nossa preocupação era \emph{forma funcional}: a relação entre $X$ e $Y$
pode não ser linear. Métodos não paramétricos resolvem exatamente esse problema ---
estimam $g(X)$ de forma flexível, sem impor uma reta.

Mas existe outro problema central em econometria aplicada, completamente distinto do
problema de forma funcional: \textbf{endogeneidade}.

A endogeneidade aparece quando $X$ está correlacionado com fatores não observados que
também afetam $Y$. Em termos do modelo \eqref{eq:linear}, isso significa que
$\mathrm{Cov}(X, \varepsilon) \neq 0$. Quando isso acontece, o estimador de OLS
(ou qualquer estimador que trate $X$ como exógeno) é viesado e não tem interpretação
causal.

\begin{exemplobox}[Exemplo clássico: educação e salário]
Suponha que queremos estimar o efeito causal da educação sobre o salário. O problema é
que pessoas com mais educação tendem a ter mais habilidade não observada, mais motivação,
melhores conexões familiares e acesso a redes de contato mais poderosas. Essas
características também afetam salários diretamente. Portanto, quando comparamos pessoas
com mais e menos educação, parte da diferença salarial que observamos não é causada pela
educação em si, mas por essas outras características. O coeficiente de OLS mistura o
efeito causal da educação com a seleção de quem decide estudar mais.
\end{exemplobox}

Aqui está o ponto crucial desta seção: \textbf{um modelo mais flexível não resolve
a endogeneidade}. Se estimarmos $g(X)$ por local linear regression sem lidar com o
fato de que $X$ é correlacionado com fatores não observados, estaremos estimando
a relação observacional entre $X$ e $Y$, não o efeito causal de $X$ sobre $Y$.

A função $g(X)$ estimada de forma não paramétrica pode revelar uma relação bela e
complexa nos dados, mas continuará sendo uma descrição da correlação, não da
causalidade. Flexibilidade de forma funcional e identificação causal são problemas
distintos, que exigem soluções distintas.

% ============================================================
\chapter{O caso linear com endogeneidade: 2SLS e control function}
% ============================================================

Vamos começar pelo caso mais simples: o modelo linear com endogeneidade.

O modelo estrutural é:

\begin{equation}
Y = \alpha + \beta X + \varepsilon,
\label{eq:structural}
\end{equation}

\noindent
mas $\mathrm{Cov}(X, \varepsilon) \neq 0$ --- $X$ está correlacionado com o erro.

Para resolver esse problema, precisamos de um \textbf{instrumento} $Z$: uma variável
que afeta $X$, mas que não tem efeito direto sobre $Y$ (exceto por meio de $X$), e que
não está correlacionada com o termo de erro $\varepsilon$.

\begin{destacebox}[O que faz um bom instrumento?]
O instrumento $Z$ precisa satisfazer duas condições: \textbf{relevância} ($Z$ precisa
ter relação com $X$) e \textbf{exogeneidade} ($Z$ não pode estar correlacionado com
$\varepsilon$). No exemplo de educação e salário, um instrumento clássico é a
\emph{distância até a faculdade mais próxima}: ela afeta a decisão de ir à faculdade,
mas não deveria afetar diretamente o salário futuro --- exceto pela via da educação.
\end{destacebox}

\section*{2SLS: usar só a parte exógena de $X$}

O estimador de \textbf{variáveis instrumentais de dois estágios} (2SLS) funciona assim:

\textbf{Primeiro estágio:} Regredimos $X$ em $Z$ (e quaisquer controles adicionais):
\begin{equation}
X = \gamma_0 + \gamma_1 Z + \nu.
\label{eq:first_stage}
\end{equation}
Obtemos a parte prevista de $X$: $\hat{X} = \hat{\gamma}_0 + \hat{\gamma}_1 Z$.

\textbf{Segundo estágio:} Regredimos $Y$ em $\hat{X}$ (em vez de $X$):
\begin{equation}
Y = \alpha + \beta \hat{X} + \text{erro}.
\end{equation}

A intuição é direta: $\hat{X}$ é a parte de $X$ que é explicada pelo instrumento $Z$.
Como $Z$ é exógeno (não correlacionado com $\varepsilon$), $\hat{X}$ também é exógeno.
Ao usar $\hat{X}$ no lugar de $X$, ``limpamos'' $X$ da parte que está contaminada
pelos fatores não observados.

\section*{Control function: controlar a parte endógena}

A \textbf{abordagem de control function} chega ao mesmo resultado por um caminho
diferente.

\textbf{Primeiro estágio:} Mesma regressão de \eqref{eq:first_stage}. Agora, guardamos
o \textbf{resíduo}:
\begin{equation}
\hat{\nu}_i = X_i - \hat{\gamma}_0 - \hat{\gamma}_1 Z_i.
\end{equation}

\textbf{Segundo estágio:} Regredimos $Y$ em $X$ \emph{e} $\hat{\nu}$:
\begin{equation}
Y = \alpha + \beta X + \delta \hat{\nu} + \text{erro}.
\label{eq:cf_linear}
\end{equation}

Por que isso funciona? O resíduo $\hat{\nu}$ captura exatamente a parte de $X$ que
não é explicada por $Z$ --- ou seja, a parte potencialmente correlacionada com os
fatores não observados em $\varepsilon$. Ao incluir $\hat{\nu}$ na equação, o
pesquisador está, em essência, controlando diretamente a ``contaminação'' de $X$.
Com $\hat{\nu}$ incluído, a variação remanescente em $X$ é, por construção, ortogonal
ao erro --- e o coeficiente $\beta$ pode ser interpretado causalmente.

\begin{destacebox}[2SLS e control function no modelo linear: resultados equivalentes]
No modelo linear clássico, 2SLS e control function produzem \emph{exatamente} o mesmo
estimador para $\beta$. Eles são duas formas diferentes de pensar sobre o mesmo
procedimento. O 2SLS olha pela janela da \emph{variação exógena}: usa apenas a parte
de $X$ que vem do instrumento. A control function olha pela janela do \emph{controle
direto}: mantém $X$ original, mas controla o componente endógeno por meio do resíduo.
No ambiente linear, ambas as janelas mostram o mesmo cenário.
\end{destacebox}

% ============================================================
\chapter{Por que o 2SLS deixa de ser suficiente em modelos não lineares}
% ============================================================

A equivalência entre 2SLS e control function depende de uma condição fundamental:
a \textbf{linearidade do modelo}. Quando relaxamos essa condição, os dois métodos
divergem — e o 2SLS começa a apresentar limitações sérias.

\section*{O que muda quando o modelo é não linear?}

No modelo linear, o efeito de $X$ sobre $Y$ é resumido por um único número, $\beta$.
Esse coeficiente tem uma interpretação clara: é o efeito médio de $X$ sobre $Y$ para
a população coberta pelo instrumento.

Mas, quando a relação verdadeira é não linear, a ideia de ``um único coeficiente
resume tudo'' começa a se desfazer. Considere os seguintes cenários:

\begin{itemize}
  \item O efeito da educação sobre salário pode ser muito diferente para quem tem
  apenas ensino fundamental do que para quem já tem ensino superior.
  \item O efeito de uma transferência de renda pode ser muito maior para famílias
  abaixo da linha de pobreza do que para famílias já acima dela.
  \item O efeito de uma regulação ambiental sobre produtividade pode depender de
  características não observadas das firmas, como sua capacidade tecnológica.
\end{itemize}

Em todos esses casos, o pesquisador pode estar interessado na \textbf{função causal
inteira} $g(X)$ — como o efeito de $X$ varia ao longo de seu suporte — e não apenas
num efeito médio.

\section*{O problema de substituir $X$ por $\hat{X}$ em modelos não lineares}

No 2SLS, a ideia central é substituir $X$ por $\hat{X}$ no segundo estágio. No
modelo linear, isso é exato. Mas imagine que o modelo verdadeiro é:

\begin{equation}
Y = g(X) + \varepsilon,
\end{equation}

\noindent
com $X$ endógeno. Se simplesmente substituirmos $X$ por $\hat{X}$ e estimarmos
$g(\hat{X})$, teremos um problema: $\hat{X}$ é uma variável contínua cuja distribuição
é diferente da de $X$. Ao compor a função não linear $g$ com $\hat{X}$, estamos
avaliando $g$ em pontos que não correspondem aos pontos originais de interesse. O
resultado é, em geral, inconsistente para $g(X)$.

\begin{destacebox}[Intuição do problema]
Imagine que $g(X) = X^2$. Se usarmos $\hat{X}$ no lugar de $X$:
$g(\hat{X}) = \hat{X}^2 = (X - \hat{\nu})^2 = X^2 - 2X\hat{\nu} + \hat{\nu}^2.$
Isso é claramente diferente de $g(X) = X^2$. A não-linearidade faz com que a
substituição mecânica de $X$ por $\hat{X}$ gere termos adicionais indesejados. O
2SLS, como ferramenta naturalmente linear, não resolve o problema corretamente.
\end{destacebox}

Essa limitação do 2SLS em ambientes não lineares motiva a necessidade de uma
abordagem diferente — e é exatamente aqui que a control function se torna mais
natural e mais poderosa.

% ============================================================
\chapter{Control function em modelos não lineares e não paramétricos}
% ============================================================

A grande vantagem da control function em relação ao 2SLS é que ela não depende da
linearidade do modelo. Sua lógica se generaliza naturalmente para ambientes mais
complexos.

\section*{A ideia central}

Recorde o que acontecia no caso linear. No segundo estágio, incluíamos $X$ e o
resíduo $\hat{\nu}$ do primeiro estágio. O resíduo capturava a parte endógena de $X$.
Ao controlar por $\hat{\nu}$, tornávamos $X$ ``como se fosse exógeno'' condicionalmente.

A mesma lógica se aplica quando a relação entre $Y$ e $X$ é não linear. O modelo
estrutural geral é:

\begin{equation}
Y = g(X, \nu) + \text{erro},
\end{equation}

\noindent
onde $\nu$ representa os fatores não observados que tornam $X$ endógeno (por exemplo,
a habilidade não observada que afeta tanto a escolha de educação quanto o salário).
O problema é que $\nu$ não é observado — mas podemos \emph{estimar} $\nu$ usando
o primeiro estágio com o instrumento $Z$.

\textbf{Primeiro estágio:} Modelamos a escolha de $X$ como função de $Z$ e de
características observadas. O resíduo do primeiro estágio, $\hat{\nu}$, é a nossa
estimativa do componente não observado que torna $X$ endógeno.

\textbf{Segundo estágio:} Estimamos a relação entre $Y$ e $X$, controlando por
$\hat{\nu}$. No modelo não linear, isso pode ser feito de maneira flexível:

\begin{equation}
Y = g(X, \hat{\nu}) + \text{erro},
\end{equation}

\noindent
onde $g(\cdot, \cdot)$ pode ser estimada não parametricamente. Com $\hat{\nu}$
incluído como controle, a variação em $X$ que é usada para identificar $g$ é, por
construção, independente do componente endógeno.

\section*{Modelos não separáveis e heterogeneidade de efeitos}

A literatura mais avançada, a partir dos trabalhos seminais de
\citet{imbens_identification_2009} e \citet{chesher_identification_2003}, mostra que
a abordagem de control function pode ser usada para identificar efeitos causais em
modelos \textbf{não separáveis}:

\begin{equation}
Y = g(X, U),
\end{equation}

\noindent
onde $U$ representa heterogeneidade individual não observada e $g$ pode ser uma
função geral (não necessariamente aditiva em $U$). Essa classe de modelos é muito
mais rica que o modelo linear, pois o efeito de $X$ sobre $Y$ pode variar de acordo
com características não observadas dos indivíduos.

Alguns conceitos úteis para entender essa literatura:

\begin{itemize}
  \item \textbf{Modelo não separável}: o efeito de $X$ sobre $Y$ pode depender de
  características não observadas $U$. Em vez de $Y = g(X) + \varepsilon$, temos
  $Y = g(X, U)$. Isso permite que o efeito causal seja diferente para diferentes
  tipos de indivíduos.

  \item \textbf{Heterogeneidade de efeitos}: o efeito causal de $X$ sobre $Y$ não é
  necessariamente o mesmo para todos os indivíduos. Pode ser maior para alguns grupos,
  menor para outros.

  \item \textbf{Variável de controle}: no contexto da control function, é a variável
  construída a partir do primeiro estágio que, quando incluída como controle, torna
  $X$ exógeno condicionalmente. No caso linear, é simplesmente o resíduo do primeiro
  estágio. Em casos mais gerais, pode ser uma transformação desse resíduo.

  \item \textbf{Identificação}: o conjunto de hipóteses que permite interpretar a
  relação estimada como causal. Na abordagem de control function em modelos não
  separáveis, as hipóteses tipicamente envolvem a forma como o instrumento afeta
  $X$ e como os fatores não observados entram no modelo.
\end{itemize}

\begin{exemplobox}[O exemplo de educação e salário revisitado]
Suponha que diferentes pessoas têm diferentes habilidades não observadas $U$, e que
tanto a decisão de estudar mais (escolha de $X$) quanto o salário ($Y$) dependem de
$U$. Em um modelo não separável, o retorno da educação pode ser diferente para pessoas
com alta habilidade e para pessoas com baixa habilidade.

A control function, nesse contexto, usa o instrumento $Z$ (por exemplo, distância à
faculdade) para construir uma estimativa de $U$ para cada indivíduo. Com essa
estimativa como controle, podemos estimar como o efeito da educação sobre o salário
varia ao longo da distribuição de habilidade --- algo impossível com um 2SLS linear
padrão.
\end{exemplobox}

\textbf{A mensagem central desta seção}: quando há endogeneidade e queremos estimar
relações causais flexíveis, a control function é uma ponte natural entre os instrumentos
e os métodos não paramétricos. Ela permite combinar a lógica de identificação dos
instrumentos com a flexibilidade dos modelos não paramétricos, sem as restrições que
o 2SLS impõe ao ambiente linear.

% ============================================================
\chapter{Onde entram os métodos semiparamétricos?}
% ============================================================

Até agora, percorremos dois extremos: o modelo linear, muito estruturado, e os modelos
não paramétricos, muito flexíveis. Na prática, muitas aplicações vivem no meio-termo.

\textbf{Métodos paramétricos} (como OLS) impõem muita estrutura: a forma funcional é
completamente especificada antes de ver os dados. Isso facilita a estimação e a
interpretação, mas ao custo de viés se a especificação estiver errada.

\textbf{Métodos não paramétricos} impõem pouca estrutura: deixam os dados revelarem
a função $g(X)$. Isso é ótimo para exploração, mas exige grandes amostras, sofre com
muitas covariáveis e pode dificultar a interpretação.

\textbf{Métodos semiparamétricos} ficam no meio: mantêm alguma parte do modelo com
estrutura paramétrica, mas deixam outra parte flexível. A ideia é combinar o melhor
dos dois mundos: a interpretabilidade dos modelos paramétricos com a flexibilidade dos
não paramétricos.

Por que isso é útil? Em muitas situações práticas, o pesquisador:

\begin{itemize}
  \item Tem uma quantidade moderada de dados — suficiente para um modelo paramétrico,
  mas talvez insuficiente para um não paramétrico completamente livre.
  \item Quer preservar alguma interpretação econômica clara para parte do modelo
  (por exemplo, o efeito de uma variável de política específica).
  \item Sabe que a relação entre certas variáveis é razoavelmente bem descrita por
  uma forma linear, mas não quer impor linearidade em toda a equação.
\end{itemize}

Exemplos de modelos semiparamétricos:

\begin{itemize}
  \item \textbf{Modelos parcialmente lineares}: $Y = \beta X_1 + g(X_2) + \varepsilon$.
  O efeito de $X_1$ é linear e tem interpretação direta; o efeito de $X_2$ é estimado
  de forma não paramétrica.
  \item \textbf{Modelos aditivos generalizados (GAM)}: $Y = \sum_j g_j(X_j) + \varepsilon$.
  Cada covariável tem seu próprio componente não paramétrico, mas esses componentes são
  aditivos, o que evita a maldição da dimensionalidade.
  \item \textbf{Regressão quantílica}: impõe estrutura linear para o quantil condicional
  de $Y$, mas não especifica a distribuição completa dos erros.
\end{itemize}

% ============================================================
\chapter{Regressão quantílica: olhando além da média}
% ============================================================

A regressão linear tradicional estima o \textbf{efeito médio} de $X$ sobre $Y$. Mais
precisamente, ela modela a \emph{média condicional} de $Y$ dado $X$:
$\mathbb{E}[Y \mid X]$.

Mas muitas perguntas aplicadas não são sobre médias:

\begin{itemize}
  \item A educação aumenta os salários igualmente para trabalhadores no percentil 10
  (salários baixos) e no percentil 90 (salários altos)?
  \item Uma política de transferência de renda melhora mais as famílias na base da
  distribuição ou as famílias que já estão perto da linha de pobreza?
  \item Uma intervenção reduz desigualdade ou apenas eleva a média sem mudar a dispersão?
\end{itemize}

Para responder a essas perguntas, precisamos estimar o efeito de $X$ em
\emph{diferentes partes da distribuição condicional de $Y$} --- não apenas na média.
É exatamente isso que a \textbf{regressão quantílica} \citep{koenker_regression_1978}
permite.

\section*{A ideia básica}

Para um nível quantílico $\tau \in (0,1)$, a regressão quantílica modela o
$\tau$-ésimo quantil condicional de $Y$:

\begin{equation}
Q_Y(\tau \mid X) = \alpha(\tau) + \beta(\tau) X.
\label{eq:qr}
\end{equation}

\noindent
Os coeficientes $\alpha(\tau)$ e $\beta(\tau)$ \emph{dependem de $\tau$}: eles podem
(e geralmente vão) ser diferentes para a mediana ($\tau = 0{,}5$) e para o percentil
90 ($\tau = 0{,}9$).

Isso é muito mais informativo do que um único coeficiente médio. Por exemplo:

\begin{itemize}
  \item Se $\beta(0{,}1) < \beta(0{,}9)$, o efeito de $X$ sobre $Y$ é maior para
  quem já está em posição mais favorável na distribuição — o que sugere que $X$
  aumenta a desigualdade.
  \item Se $\beta(0{,}1) > \beta(0{,}9)$, o efeito é maior para quem está em posição
  pior — o que sugere que $X$ reduz a desigualdade.
  \item Se $\beta(\tau)$ é aproximadamente constante em $\tau$, o efeito é homogêneo
  ao longo da distribuição.
\end{itemize}

\begin{exemplobox}[Exemplo: efeito da educação ao longo da distribuição salarial]
\citet{angrist_quantile_2006} e outros autores mostram que o retorno da educação é
heterogêneo: para trabalhadores no topo da distribuição salarial, os retornos adicionais
de um ano de escolaridade são maiores do que para trabalhadores na base. Uma regressão
linear simples calcularia uma média desses efeitos, mascarando essa heterogeneidade
importante. A regressão quantílica revela o padrão completo.
\end{exemplobox}

\section*{Por que é semiparamétrica?}

A regressão quantílica é classificada como semiparamétrica porque:

\begin{itemize}
  \item Ela impõe uma estrutura \emph{parcial}: assume que o quantil condicional é
  linear em $X$ (equação \ref{eq:qr}).
  \item Mas \emph{não} especifica a distribuição completa dos erros. Ao contrário da
  regressão linear (que sob hipóteses padrão assume normalidade ou pelo menos
  homocedasticidade), a regressão quantílica é robusta a distribuições assimétricas,
  caudas pesadas e heterocedasticidade.
\end{itemize}

\begin{notatecnica}[Nota técnica: como a regressão quantílica é estimada?]
A regressão quantílica minimiza uma função de perda assimétrica, chamada \emph{função
check} (ou perda tilted $\ell_1$):
\[
\hat{\beta}(\tau) = \arg\min_b \sum_{i=1}^n \rho_\tau(Y_i - X_i' b),
\]
onde $\rho_\tau(u) = u[\tau - \mathbf{1}(u < 0)]$. Para $\tau = 0{,}5$, isso equivale
à minimização do erro absoluto médio (regressão de mediana). Para outros valores de
$\tau$, a função penaliza erros positivos e negativos de forma diferente.
\end{notatecnica}

\section*{Limitações importantes}

Assim como outros métodos, a regressão quantílica tem limitações que merecem atenção:

\begin{itemize}
  \item \textbf{Endogeneidade}: a regressão quantílica padrão não resolve o problema
  de endogeneidade. Se $X$ é endógeno, os quantis estimados continuarão confundindo
  causalidade com seleção. Há extensões da regressão quantílica para variáveis
  instrumentais (como em \citealt{chernozhukov_instrumental_2005}), mas elas são
  mais complexas e exigem cuidado na interpretação.
  \item \textbf{Interpretação dos quantis condicionais}: os quantis condicionais
  estimados descrevem a distribuição condicional de $Y$ dado $X$, não
  necessariamente o efeito estrutural para indivíduos num determinado quantil da
  distribuição de $Y$. A distinção entre quantis condicionais e efeitos de tratamento
  quantílicos é sutil, mas importante.
\end{itemize}

% ============================================================
\chapter{Como escolher entre os métodos?}
% ============================================================

Chegamos a um ponto em que temos várias ferramentas disponíveis. Esta seção oferece
um guia prático para a escolha entre elas, na forma de uma sequência de perguntas.

\section*{Perguntas de diagnóstico}

\textbf{1. A relação entre $X$ e $Y$ parece aproximadamente linear?}

\begin{itemize}
  \item Se \textbf{sim}: a regressão linear é um bom ponto de partida. Ela é simples,
  interpretável e eficiente quando a linearidade é uma aproximação razoável.
  \item Se \textbf{não} (ou se não sabe): métodos não paramétricos, como local linear
  regression, podem revelar padrões que uma reta esconde. Comece explorando
  graficamente a relação antes de escolher um modelo.
\end{itemize}

\textbf{2. O objetivo é descrição/previsão ou identificação causal?}

\begin{itemize}
  \item Se o objetivo é \textbf{descrever padrões} nos dados ou \textbf{prever} $Y$
  a partir de $X$: métodos não paramétricos são ferramentas poderosas.
  \item Se o objetivo é \textbf{identificar um efeito causal}: a forma funcional
  importa menos do que a estratégia de identificação. É necessário pensar em
  exogeneidade, seleção e outras ameaças à validade interna.
\end{itemize}

\textbf{3. Existe endogeneidade?}

\begin{itemize}
  \item Se \textbf{sim}: é necessário um instrumento, uma estratégia de
  diferenças-em-diferenças, um desenho de regressão descontínua, ou outra fonte de
  variação exógena. A escolha do método de estimação (2SLS, control function) vem
  \emph{depois} da estratégia de identificação.
  \item Se \textbf{não} (ou se a exogeneidade for defensável): OLS e métodos
  não paramétricos são viáveis.
\end{itemize}

\textbf{4. O modelo causal é bem representado por um efeito linear médio?}

\begin{itemize}
  \item Se \textbf{sim}: o 2SLS é adequado e bem-fundamentado.
  \item Se \textbf{não} (relação não linear, heterogeneidade, modelo não separável):
  a abordagem de control function, possivelmente combinada com estimação não
  paramétrica no segundo estágio, é mais apropriada.
\end{itemize}

\textbf{5. Você quer estudar efeitos em diferentes partes da distribuição de $Y$?}

\begin{itemize}
  \item Se \textbf{sim}: a regressão quantílica é a ferramenta natural. Ela permite
  estudar como o efeito de $X$ varia ao longo da distribuição condicional de $Y$.
\end{itemize}

\section*{Limitações que merecem atenção}

Nenhum método é livre de limitações. A tabela abaixo resume as principais:

\medskip
\begin{center}
\begin{tabular}{ll}
\hline
\textbf{Método} & \textbf{Limitações principais} \\
\hline
Regressão linear & Viés se a relação não for linear \\
Local linear regression & Precisa de muitos dados; sofre com muitas covariáveis \\
& Escolha do bandwidth pode ser determinante \\
2SLS & Naturalmente linear; perde consistência em modelos não lineares \\
Control function & Exige hipóteses fortes; sensível à especificação do 1.º estágio \\
& Instrumento precisa ser válido \\
Regressão quantílica & Não resolve endogeneidade automaticamente \\
& Interpretação de quantis condicionais requer cuidado \\
\hline
\end{tabular}
\end{center}

% ============================================================
\chapter{Conclusão}
% ============================================================

Este guia percorreu um caminho que começa na regressão linear e termina nos métodos
semiparamétricos, passando pelos métodos não paramétricos e pelo problema de
endogeneidade.

A \textbf{regressão linear} é um ponto de partida poderoso e indispensável em
econometria aplicada. Mas ela impõe duas simplificações importantes: linearidade
(o efeito de $X$ sobre $Y$ é constante) e foco no efeito médio (um único número
resume a relação).

\textbf{Métodos não paramétricos}, como a regressão linear local, relaxam a forma
funcional e permitem enxergar relações mais flexíveis nos dados. Eles são ferramentas
valiosas de exploração e, quando as hipóteses de identificação são satisfeitas, também
de estimação causal. Mas flexibilidade na forma funcional não elimina o problema de
endogeneidade.

Para \textbf{identificação causal}, é necessário uma estratégia específica. No caso
linear, o \textbf{2SLS} e a \textbf{control function} oferecem soluções equivalentes,
cada uma com sua intuição distinta. Em modelos não lineares, não separáveis ou com
heterogeneidade de efeitos, a control function se torna mais natural: ela generaliza
diretamente para esses ambientes, enquanto o 2SLS fica preso à lógica linear.

\textbf{Métodos semiparamétricos}, como a \textbf{regressão quantílica}, aparecem
como um meio-termo útil: mais flexíveis do que modelos lineares simples (pois não se
limitam a um efeito médio), mas ainda interpretáveis e viáveis em aplicações com
amostras de tamanho moderado.

Conhecer esses métodos não é apenas ampliar a caixa de ferramentas. É desenvolver
a capacidade de fazer as perguntas certas: a relação é linear? O efeito é homogêneo?
A estimação é causal? As respostas a essas perguntas é que guiam a escolha do método
— e não o contrário.

\bigskip

\begin{destacebox}[Mensagem final]
Métodos são meios, não fins. O passo mais importante antes de escolher qualquer
ferramenta é entender a pergunta que está sendo feita e as ameaças à sua resposta.
Um pesquisador que entende profundamente a regressão linear e seus limites está
mais bem equipado do que um que conhece muitos métodos sem entender suas hipóteses.
\end{destacebox}

% ============================================================
%  Referências para leitura posterior
% ============================================================
\chapter*{Leituras recomendadas}
\addcontentsline{toc}{chapter}{Leituras recomendadas}

As referências abaixo são um ponto de partida para quem quiser aprofundar os tópicos
deste guia.

\textbf{Sobre regressão não paramétrica e estimação local:}

\begin{itemize}
  \item \citet{fan_local_1996} — referência clássica sobre regressão polinomial
  local; rigorosa, mas acessível a quem tem boa base de estatística.
  \item \citet{li_nonparametric_2007} — manual abrangente de econometria não
  paramétrica e semiparamétrica, com muitas aplicações.
  \item \citet{wand_kernel_1994} — texto acessível sobre suavização por kernel,
  incluindo escolha de bandwidth.
\end{itemize}

\textbf{Sobre control function em modelos não lineares e não separáveis:}

\begin{itemize}
  \item \citet{imbens_identification_2009} — artigo fundamental sobre identificação
  e estimação em modelos não separáveis com variáveis instrumentais; usa a abordagem
  de control function de forma geral.
  \item \citet{chesher_identification_2003} — mostra condições de identificação em
  modelos não separáveis com heterogeneidade de efeitos; mais técnico, mas essencial
  para quem quiser entender os fundamentos.
\end{itemize}

\textbf{Sobre regressão quantílica:}

\begin{itemize}
  \item \citet{koenker_regression_1978} — artigo original de Koenker e Bassett que
  introduz a regressão quantílica.
  \item \citet{koenker_quantile_2005} — livro completo de Koenker sobre regressão
  quantílica, teoria e aplicações.
  \item \citet{chernozhukov_instrumental_2005} — extensão da regressão quantílica para
  variáveis instrumentais.
\end{itemize}

\textbf{Introduções gerais a econometria não paramétrica:}

\begin{itemize}
  \item \citet{pagan_nonparametric_1999} — introdução acessível a métodos não
  paramétricos para economistas; bom ponto de partida para quem está começando.
  \item \citet{cattaneo_simple_2020} — apresenta estimadores de regressão não
  paramétrica com foco em aplicações práticas e inferência, com exemplos em R e Stata.
\end{itemize}
